\chapter{Infinitely Divisible and Stable Laws}
\begin{center}\itshape
Which laws can be the limit of a sum of many small independent pieces? The Gaussian, of course; the Poisson, when the pieces are rare events;\\
and a whole family in between, described by one formula. Among them are the stable laws,\\
the only possible limits of rescaled sums of i.i.d.\ variables, and the right model for data whose averages refuse to settle down.
\end{center}

In 1824 Sim\'eon-Denis Poisson, studying how to combine astronomical observations, noticed something disturbing. If the errors have the density
$\frac1{\pi(1+x^2)}$, the average of any number of observations has exactly the same law as a single one: averaging gains nothing. Nobody paid much
attention, Poisson included. In 1853 Augustin Cauchy came across the same density during a public dispute with Ir\'en\'ee-Jules Bienaym\'e about the method of
least squares, and this time it stuck, with his name on it. The curve itself was older still: it is the ``witch of Agnesi'' of eighteenth-century
geometry. Stephen Stigler told this story in 1974.

Volume~3 met the Cauchy law as a curiosity. Paul L\'evy, in the 1920s, saw it as the first member of a family. He asked which laws $\mu$ have the property
that a sum of independent copies, rescaled, has law $\mu$ again; for the Gaussian the scale is $\sqrt n$, for Cauchy it is $n$, and L\'evy found all the
others, with scales $n^{1/\alpha}$ for $0<\alpha\le2$, in his book of 1925. Then Bruno de Finetti asked in 1929 a question about processes with
independent increments that leads to a bigger family: the laws that can be split, for every $n$, into $n$ independent identical pieces. Kolmogorov
described them in 1932 when the variance is finite, L\'evy in general in 1934, and Aleksandr Khintchine gave the proof that textbooks still follow in 1937.
The result, the \emph{L\'evy--Khintchine formula}, says that every such law is a Gaussian plus a superposition of Poisson clouds of jumps.

This chapter proves the formula, shows that it describes the limits of sums of small independent pieces, classifies the symmetric stable laws and
their domains of attraction, and ends with a stable law that has been hiding in Volume~3 all along: the time a fair gambler waits to get ahead.

\begin{reminder}[from Volume 3 and Chapters 3 and 4]
L\'evy's continuity theorem (Theorem~\ref{thm:levycont}): if characteristic functions $\varphi_n$ converge pointwise to a function $\varphi$ continuous at $0$,
then $\varphi$ is a characteristic function and the laws converge. A \emph{compound Poisson} variable is $\sum_{i\le N}\xi_i$ with $N$ Poisson$(\lambda)$
and $\xi_i$ i.i.d.\ with law $\nu$, independent of $N$; conditioning on $N$, its characteristic function is
$\sum_k e^{-\lambda}\frac{\lambda^k}{k!}\varphi_\nu(t)^k=\exp\bigl(\lambda(\varphi_\nu(t)-1)\bigr)$. The Cauchy law has characteristic function $e^{-|t|}$
(Volume~3, Corollary~6.9). For a Brownian motion, the hitting time $T_a$ of $a>0$ satisfies $P(T_a\le t)=2(1-\Phi(a/\sqrt t))$ and $\mathbb ET_a=\infty$
(Volume~3, Corollary~11.13), and the process starts afresh at stopping times (Theorem~\ref{thm:strongmarkov}). Finally, for complex $w$ with $|w|\le\frac12$,
$|\log(1+w)-w|\le|w|^2$ (proof of Theorem~\ref{thm:feller}).
\end{reminder}

\section{Infinitely divisible laws}

A law $\mu$ on $\R$ (or a random variable, or its characteristic function $\varphi$) is \emph{infinitely divisible} if for every $n\ge1$ there is a law $\mu_n$
with $\mu=\mu_n^{*n}$, that is, $\varphi=\varphi_n^n$ for a characteristic function $\varphi_n$: a variable with law $\mu$ is a sum of $n$ i.i.d.\ pieces.
Examples: $N(m,\sigma^2)$ (pieces $N(m/n,\sigma^2/n)$); Poisson$(\lambda)$ (pieces Poisson$(\lambda/n)$); any compound Poisson law (replace $\lambda$ by $\lambda/n$);
Cauchy ($e^{-|t|}=(e^{-|t|/n})^n$); the gamma laws (Exercise~\ref{ex:idexamples}). Non-examples: any non-constant bounded variable
(Exercise~\ref{ex:bounded}), so the uniform and the Bernoulli laws are not infinitely divisible.

To take $n$-th roots of characteristic functions we need logarithms.

\begin{lemma}[continuous logarithm]\label{lem:contlog}
If $\varphi:\R\to\C$ is continuous, never $0$, and $\varphi(0)=1$, there is a unique continuous $L:\R\to\C$ with $L(0)=0$ and $e^L=\varphi$.
\end{lemma}
\begin{proof}
Uniqueness: two such functions differ by a continuous function with values in $2\pi i\Z$, hence constant, and equal to $0$ at $0$. Existence: fix $T>0$. On
$[-T,T]$, $|\varphi|\ge\delta$ for some $\delta>0$, and $\varphi$ is uniformly continuous, so there is $\eta>0$ with $|\varphi(s)-\varphi(s')|<\delta$ when
$|s-s'|\le\eta$; then $|\varphi(s)/\varphi(s')-1|<1$ and the principal logarithm $\operatorname{Log}$ of the ratio is defined and continuous. For $t\in[k\eta,(k+1)\eta]$,
$0\le k\eta<T$, put
\[
L(t)=\sum_{j<k}\operatorname{Log}\frac{\varphi((j+1)\eta)}{\varphi(j\eta)}+\operatorname{Log}\frac{\varphi(t)}{\varphi(k\eta)},
\]
and symmetrically for $t<0$. The definitions agree at the grid points, $L$ is continuous, and the product telescopes to $e^{L(t)}=\varphi(t)$. By uniqueness
the functions built for different $T$ agree, so they define $L$ on $\R$.
\end{proof}

\begin{proposition}\label{prop:idroot}
An infinitely divisible characteristic function $\varphi$ never vanishes. If $L$ is its continuous logarithm, the only characteristic function $\varphi_n$ with
$\varphi_n^n=\varphi$ is $e^{L/n}$.
\end{proposition}
\begin{proof}
Let $\varphi_n$ be a characteristic function with $\varphi_n^n=\varphi$. Then $|\varphi_n|^2$ is a characteristic function (of $Y-Y'$, with $Y,Y'$ independent with
characteristic function $\varphi_n$), equal to $|\varphi|^{2/n}$. As $n\to\infty$, $|\varphi(t)|^{2/n}$ tends to $g(t)=1$ if $\varphi(t)\ne0$ and to $0$ otherwise. As
$\varphi$ is continuous with $\varphi(0)=1$, $g=1$ near $0$, so $g$ is continuous at $0$, and by L\'evy's theorem it is a characteristic function, hence continuous
everywhere. A continuous function with values in $\{0,1\}$ and $g(0)=1$ is identically $1$: $\varphi$ never vanishes. Then $\varphi_n$ never vanishes either,
and if $L_n$ is its continuous logarithm, $nL_n$ and $L$ are both continuous logarithms of $\varphi$, so $nL_n=L$.
\end{proof}

\begin{think}
If $X$ is infinitely divisible, is $X$ plus an independent constant? Is $aX$? Is $X+Y$, with $Y$ independent and infinitely divisible? And if a sequence of
infinitely divisible laws converges, is the limit infinitely divisible? (Use Proposition~\ref{prop:idroot}: the $n$-th roots converge too.)
\end{think}

\section{The L\'evy--Khintchine formula}

A \emph{L\'evy measure} is a measure $\Pi$ on $\R\setminus\{0\}$ with $\int\min(1,x^2)\,\Pi(\dd x)<\infty$: it may have infinite mass, but only near $0$, and
not too much. For $b\in\R$, $\sigma^2\ge0$ and a L\'evy measure $\Pi$, put
\begin{equation}\label{eq:lk}
\psi(t)=ibt-\frac{\sigma^2t^2}2+\int\bigl(e^{itx}-1-itx\mathbf 1_{|x|\le1}\bigr)\,\Pi(\dd x).
\end{equation}
The integral converges: by~\eqref{eq:taylor}, the integrand is at most $\frac{t^2x^2}2$ for $|x|\le1$ and at most $2$ for $|x|>1$. By dominated
convergence, $\psi$ is continuous, and $\psi(0)=0$. Read~\eqref{eq:lk} as a recipe: a drift $b$, a Gaussian part $\sigma^2$, and jumps of size $x$ arriving
at rate $\Pi(\dd x)$, with the small jumps compensated by their mean so that infinitely many of them can be added up.

\begin{theorem}[L\'evy--Khintchine]\label{thm:lk}
\begin{enumerate}[label=(\alph*)]
\item For every triple $(b,\sigma^2,\Pi)$, $e^\psi$ is an infinitely divisible characteristic function.
\item Every infinitely divisible characteristic function is $e^\psi$ for some triple.
\item The triple is determined by the law.
\end{enumerate}
\end{theorem}
\begin{proof}
(a) For $\varepsilon\in(0,1)$, the restriction $\Pi_\varepsilon$ of $\Pi$ to $\{|x|>\varepsilon\}$ has finite mass $\lambda_\varepsilon$, and $m_\varepsilon=\int_{\varepsilon<|x|\le1}x\,\Pi(\dd x)$ is
finite. If $\lambda_\varepsilon>0$, $\exp\bigl(\int(e^{itx}-1)\Pi_\varepsilon(\dd x)\bigr)$ is the characteristic function of a compound Poisson variable with rate
$\lambda_\varepsilon$ and jump law $\Pi_\varepsilon/\lambda_\varepsilon$. So
\[
\psi_\varepsilon(t)=i(b-m_\varepsilon)t-\frac{\sigma^2t^2}2+\int_{|x|>\varepsilon}(e^{itx}-1)\,\Pi(\dd x)
\]
is the exponent of the characteristic function of an independent sum of a Gaussian and a compound Poisson variable. Since $\psi_\varepsilon(t)=ibt-\frac{\sigma^2t^2}2+\int_{|x|>\varepsilon}(e^{itx}-1-itx\mathbf 1_{|x|\le1})\Pi(\dd x)$, dominated convergence gives
$\psi_\varepsilon(t)\to\psi(t)$ as $\varepsilon\to0$, and
$e^\psi$ is continuous at $0$. By L\'evy's theorem, $e^\psi$ is a characteristic function. It is infinitely divisible because $\psi/n$ is the exponent of the triple
$(b/n,\sigma^2/n,\Pi/n)$.

(b) Let $\varphi$ be infinitely divisible, $L$ its continuous logarithm, and $\mu_n$ the law with characteristic function $\varphi_n=e^{L/n}$
(Proposition~\ref{prop:idroot}). On every compact set $L$ is bounded, so $n(\varphi_n-1)=n(e^{L/n}-1)\to L$ uniformly on compact sets. We write
$n(\varphi_n-1)$ in a form whose limits we can control. Let $c(x)=\max(-1,\min(1,x))$, a bounded continuous truncation, $w(x)=\frac{x^2}{1+x^2}$, and
\[
k(t,x)=\bigl(e^{itx}-1-itc(x)\bigr)\frac{1+x^2}{x^2}\quad(x\ne0),\qquad k(t,0)=-\frac{t^2}2.
\]
For fixed $t$, $k(t,\cdot)$ is continuous (near $0$, $c(x)=x$ and~\eqref{eq:taylor} applies) and bounded (by $2(2+|t|)$ for $|x|\ge1$). Then
\begin{equation}\label{eq:khintchine}
n(\varphi_n(t)-1)=itb_n+\int k(t,x)\,G_n(\dd x),
\end{equation}
where $b_n=n\int c\,\dd\mu_n$ and $G_n(\dd x)=n\,w(x)\,\mu_n(\dd x)$, a finite measure.

\emph{The $G_n$ are bounded and tight.} Taking real parts, $n\int(1-\cos tx)\,\mu_n(\dd x)\to R(t)=-\operatorname{Re}L(t)$, uniformly for $|t|\le1$. Averaging
over $t\in[-u,u]$, for $0<u\le1$,
\[
n\int\Bigl(1-\frac{\sin ux}{ux}\Bigr)\mu_n(\dd x)\to A(u)=\frac1{2u}\int_{-u}^uR(t)\dd t,
\]
and $A(u)\to R(0)=0$ as $u\to0$. Now $1-\frac{\sin y}y\ge c_0\frac{y^2}{1+y^2}$ for all $y$ and some $c_0>0$ (both sides are positive for $y\ne0$, and their ratio
is continuous with limits $\frac16$ at $0$ and $1$ at infinity). With $u=1$ this bounds $G_n(\R)$ by $A(1)/c_0+o(1)$. And $1-\frac{\sin y}y\ge\frac12$ for $|y|\ge2$, so
with $u=\frac2K$, $G_n(|x|>K)\le n\mu_n(|x|>K)\le2A(\frac2K)+o(1)$, which is small for $K$ large and $n\ge n_0$; increasing $K$ takes care of the finitely
many $n<n_0$.

\emph{Passing to the limit.} By Prokhorov's theorem on $\R$ (Volume~6, Theorem~11.7, applied to $G_n/G_n(\R)$ after extracting a subsequence along which the
masses converge), some subsequence $G_{n_j}$ converges weakly to a finite measure $G$. As $k(t,\cdot)$ is bounded and continuous,
$\int k(t,x)G_{n_j}(\dd x)\to\int k(t,x)G(\dd x)$. By~\eqref{eq:khintchine}, $itb_{n_j}$ then converges for every $t$, so $b_{n_j}\to b$ for some real $b$, and
\[
L(t)=itb+\int k(t,x)\,G(\dd x).
\]
Let $\sigma^2=G(\{0\})$ and $\Pi(\dd x)=\frac{1+x^2}{x^2}G(\dd x)$ on $\R\setminus\{0\}$. Then $\int\min(1,x^2)\Pi(\dd x)\le2G(\R)<\infty$, and
$L(t)=itb-\frac{\sigma^2t^2}2+\int(e^{itx}-1-itc(x))\Pi(\dd x)$. Finally $c(x)-x\mathbf 1_{|x|\le1}=\operatorname{sign}(x)\mathbf 1_{|x|>1}$ is $\Pi$-integrable, so
replacing $c$ by $x\mathbf 1_{|x|\le1}$ changes only $b$, into $b+\int_{|x|>1}\operatorname{sign}(x)\,\Pi(\dd x)$. This is~\eqref{eq:lk}.

(c) The law determines $\varphi$, hence its continuous logarithm, which is $\psi$ (it is continuous, vanishes at $0$, and $e^\psi=\varphi$). Let
\[
\theta(t)=\psi(t)-\frac12\int_{-1}^1\psi(t+s)\dd s.
\]
For the three parts of~\eqref{eq:lk}: the drift gives $ibt-\frac12\int_{-1}^1ib(t+s)\dd s=0$; the Gaussian part gives
$-\frac{\sigma^2}2\bigl(t^2-\frac12\int_{-1}^1(t+s)^2\dd s\bigr)=\frac{\sigma^2}6$; and for each $x$, the jump integrand gives $e^{itx}\bigl(1-\frac{\sin x}x\bigr)$, because the
terms $-1$ and $-itx\mathbf 1_{|x|\le1}$ are affine in $t$ and disappear like the drift. Exchanging the integrals (the integrand is bounded by a constant times
$\min(1,x^2)$ for $|s|\le1$),
\[
\theta(t)=\int e^{itx}\,\rho(\dd x),\qquad\rho=\frac{\sigma^2}6\delta_0+\Bigl(1-\frac{\sin x}x\Bigr)\Pi(\dd x).
\]
Since $0\le1-\frac{\sin x}x\le C\min(1,x^2)$, $\rho$ is a finite measure, determined by its Fourier transform $\theta$ (Volume~3, Theorem~6.8, after normalisation).
So $\psi$ determines $\rho$, hence $\sigma^2=6\rho(\{0\})$ and $\Pi$ (divide by $1-\frac{\sin x}x>0$ off $0$), and then $b$.
\end{proof}

\begin{example}\label{ex:lkexamples}
The Gaussian: $\Pi=0$. Poisson$(\lambda)$: $\Pi=\lambda\delta_1$ and $b=\lambda$, since $\lambda(e^{it}-1)=i\lambda t+\lambda(e^{it}-1-it)$. Compound Poisson with
jump law $\nu$ and rate $\lambda$: $\Pi=\lambda\nu$ (if $\nu(\{0\})=0$). Cauchy: $\Pi(\dd x)=\frac{\dd x}{\pi x^2}$ and $b=0$, because
$\int(1-\cos tx)\frac{\dd x}{\pi x^2}=\frac{|t|}\pi\int\frac{1-\cos u}{u^2}\dd u=|t|$ and the odd part of the integrand integrates to $0$. The Cauchy law is
made only of jumps, infinitely many small ones and a few very large ones.
\end{example}

\section{Limits of sums of small pieces}

The proof of Theorem~\ref{thm:lk}(b) took an infinitely divisible law apart into $n$ small pieces. Conversely, sums of small pieces converge to
infinitely divisible laws, and the triple of the limit can be read off the pieces. We treat rows of i.i.d.\ variables, which is all we need. A truncation
level $a>0$ with $\Pi(\{-a,a\})=0$ replaces $1$ in~\eqref{eq:lk}; this only changes $b$ (into $b_a=b+\int x(\mathbf 1_{|x|\le a}-\mathbf 1_{|x|\le1})\Pi(\dd x)$).

\begin{theorem}[convergence to infinitely divisible laws]\label{thm:arrays}
For each $n$ let $X_{n,1},\dots,X_{n,n}$ be i.i.d.\ with law $\mu_n$, and $S_n=\sum_kX_{n,k}$. Let $(b,\sigma^2,\Pi)$ be a triple, written with a truncation level
$a$ such that $\Pi(\{-a,a\})=0$, and assume:
\begin{enumerate}[label=(\roman*)]
\item for every $y>0$ with $\Pi(\{-y,y\})=0$,
\[
n\mu_n((y,\infty))\to\Pi((y,\infty)),\qquad n\mu_n((-\infty,-y))\to\Pi((-\infty,-y));
\]
\item $\displaystyle\lim_{\varepsilon\to0}\limsup_n\Bigl|n\int_{|x|\le\varepsilon}x^2\,\mu_n(\dd x)-\sigma^2\Bigr|=0$;
\item $n\int_{|x|\le a}x\,\mu_n(\dd x)\to b$.
\end{enumerate}
Then $S_n$ converges in law to the infinitely divisible law with exponent $\psi(t)=ibt-\frac{\sigma^2t^2}2+\int(e^{itx}-1-itx\mathbf 1_{|x|\le a})\Pi(\dd x)$.
\end{theorem}
\begin{proof}
Fix $t$, let $g(x)=e^{itx}-1-itx\mathbf 1_{|x|\le a}$, and let $0<\varepsilon<a<K$ with $\Pi(\{\pm\varepsilon\})=\Pi(\{\pm K\})=0$ (all but countably many values qualify).
Then $n(\varphi_{\mu_n}(t)-1)=itn\int_{|x|\le a}x\,\dd\mu_n+n\int g\,\dd\mu_n$, and we split the second integral into three regions.

On $|x|\le\varepsilon$: by~\eqref{eq:taylor}, $g(x)=-\frac{t^2x^2}2+r(x)$ with $|r(x)|\le\frac{|t|^3|x|^3}6\le\frac{|t|^3\varepsilon}6x^2$. So
$n\int_{|x|\le\varepsilon}g\,\dd\mu_n=-\frac{t^2}2V_n(\varepsilon)+O(\varepsilon V_n(\varepsilon))$ with $V_n(\varepsilon)=n\int_{|x|\le\varepsilon}x^2\dd\mu_n$, and on the limit side
$|\int_{|x|\le\varepsilon}g\,\dd\Pi|\le\frac{t^2}2\int_{|x|\le\varepsilon}x^2\dd\Pi\to0$ as $\varepsilon\to0$.

On $\varepsilon<|x|\le K$: by (i), the finite measures $n\mu_n$ restricted to this set have distribution functions converging to those of $\Pi$ restricted to
it, at all continuity points, and their masses converge. By Volume~3, Proposition~6.3 and the portmanteau theorem (Volume~6, Theorem~11.4), after
normalising, $n\int_{\varepsilon<|x|\le K}g\,\dd\mu_n\to\int_{\varepsilon<|x|\le K}g\,\dd\Pi$, since $g$ is bounded on this set and continuous except at $\pm a$, which
$\Pi$ does not charge.

On $|x|>K$: $|g|\le2$, and $n\mu_n(|x|>K)\to\Pi(|x|>K)$, which is small for $K$ large.

Collecting, and using (ii) and (iii), $\limsup_n|n(\varphi_{\mu_n}(t)-1)-\psi(t)|$ is bounded by a quantity that tends to $0$ when $\varepsilon\to0$ and $K\to\infty$. So
$w_n=\varphi_{\mu_n}(t)-1$ satisfies $nw_n\to\psi(t)$; in particular $w_n\to0$, and $|n\operatorname{Log}(1+w_n)-nw_n|\le n|w_n|^2=|nw_n|^2/n\to0$. Hence
$\varphi_{S_n}(t)=\exp\bigl(n\operatorname{Log}(1+w_n)\bigr)\to e^{\psi(t)}$, which is a characteristic function by Theorem~\ref{thm:lk}, and L\'evy's theorem concludes.
\end{proof}

\begin{example}[the three classical limits]\label{ex:threelimits}
(1) \emph{Central limit theorem.} If $X$ is centred with variance $\sigma^2$ and $\mu_n$ is the law of $X/\sqrt n$, then
$n\mu_n(|x|>y)\le E[X^2;|X|>y\sqrt n]/y^2\to0$, so $\Pi=0$; $V_n(\varepsilon)=\mathbb E[X^2;|X|\le\varepsilon\sqrt n]\to\sigma^2$; and
$|n\mathbb E[X/\sqrt n;|X|\le\sqrt n]|=\sqrt n\,|\mathbb E[X;|X|>\sqrt n]|\le\mathbb E[X^2;|X|>\sqrt n]\to0$. The limit is $N(0,\sigma^2)$.
(2) \emph{Law of small numbers.} If $\mu_n$ is Bernoulli$(\frac\lambda n)$, then $n\mu_n((y,\infty))=\lambda$ for $y<1$ and $0$ for $y\ge1$, so $\Pi=\lambda\delta_1$; $V_n(\varepsilon)=0$
for $\varepsilon<1$; with $a=2$, $n\int_{|x|\le2}x\,\dd\mu_n=\lambda$. The exponent is $i\lambda t+\lambda(e^{it}-1-it)=\lambda(e^{it}-1)$: the limit is Poisson$(\lambda)$.
(3) \emph{Rare jumps.} If $\mu_n=(1-\frac\lambda n)\delta_0+\frac\lambda n\nu$, the same computation gives the compound Poisson law with rate $\lambda$ and jump law $\nu$:
insurance claims, packets arriving at a router, mutations along a genome.
\end{example}

\begin{remark}
Conversely, if the sums $S_n$ of an i.i.d.\ array converge in law and $\mu_n\Rightarrow\delta_0$, then the limit is infinitely divisible and conditions (i)--(iii)
hold for its triple; the proof follows that of Theorem~\ref{thm:lk}(b), with the uniqueness in (c) forcing the whole sequence $G_n$ to converge. Together with
Theorem~\ref{thm:arrays} this describes exactly when a sum of many small independent pieces has a limit, and which; see Gnedenko and Kolmogorov's book.
\end{remark}

\section{Stable laws}

A non-degenerate symmetric random variable $X$ is \emph{stable} if for every $n\ge1$ there is $a_n>0$ such that, for i.i.d.\ copies $X_1,\dots,X_n$ of $X$,
\[
X_1+\dots+X_n\overset{\mathrm d}=a_nX.
\]
The Gaussian ($a_n=\sqrt n$) and the Cauchy law ($a_n=n$) are stable.

\begin{theorem}[symmetric stable laws]\label{thm:stable}
A symmetric variable $X$ is stable if and only if $\varphi_X(t)=e^{-c|t|^\alpha}$ for some $c>0$ and $0<\alpha\le2$. Then $a_n=n^{1/\alpha}$. The number $\alpha$ is the
\emph{index}.
\end{theorem}
\begin{proof}
\emph{Necessity.} $X\overset{\mathrm d}=(X_1+\dots+X_n)/a_n$ is a sum of $n$ i.i.d.\ pieces, so $X$ is infinitely divisible and $\varphi=\varphi_X$ never vanishes. As $X$ is
symmetric, $\varphi$ is real, continuous, never $0$ and equal to $1$ at $0$, so $0<\varphi\le1$. Let $\psi=-\log\varphi\ge0$: it is continuous, even, $\psi(0)=0$,
and $\psi\not\equiv0$ since $X$ is not degenerate. Stability says $\varphi(t)^n=\varphi(a_nt)$, that is
\begin{equation}\label{eq:stablefe}
\psi(a_nt)=n\psi(t).
\end{equation}
\emph{Step 1: $a_{mn}=a_ma_n$.} By~\eqref{eq:stablefe}, $\psi(a_{mn}t)=mn\psi(t)=\psi(a_ma_nt)$ for all $t$. If $\psi(st)=\psi(t)$ for all $t$ with $0<s<1$, then
$\psi(t)=\psi(s^kt)\to\psi(0)=0$, so $\psi\equiv0$; the case $s>1$ reduces to this one with $1/s$. So the ratio $s=a_{mn}/(a_ma_n)$ is $1$.

\emph{Step 2: $(a_n)$ is increasing.} For $m<n$, $S_n=S_m+(S_n-S_m)$ with independent symmetric terms, so $P(S_n>x)\ge P(S_m>x,\,S_n-S_m\ge0)\ge\frac12P(S_m>x)$.
Suppose $a_n<a_m$. Applying this to $n^j>m^j$ and using Step~1, $P(X>x/a_n^j)\ge\frac12P(X>x/a_m^j)$; with $x=a_m^jy$, where $P(X>y)>0$,
\[
P\Bigl(X>y\Bigl(\frac{a_m}{a_n}\Bigr)^j\Bigr)\ge\frac12P(X>y)>0\quad\text{for all }j,
\]
which is absurd since the left side tends to $0$. And $a_n=a_m$ with $n\ne m$ would give $n\psi=m\psi$ by~\eqref{eq:stablefe}. So for $m<n$ neither $a_n<a_m$ nor $a_n=a_m$ is
possible: $(a_n)$ is strictly increasing.

\emph{Step 3: $a_n=n^{1/\alpha}$.} Let $\beta=\log a_2/\log2>0$. For $n\ge2$ and $j\ge1$, pick $k$ with $2^k\le n^j<2^{k+1}$. By Steps 1 and 2,
$a_2^k\le a_n^j<a_2^{k+1}$, so $\frac kj\log a_2\le\log a_n<\frac{k+1}j\log a_2$; as $j\to\infty$, $\frac kj\to\frac{\log n}{\log2}$, hence $a_n=n^\beta$. Put $\alpha=1/\beta$.

\emph{Step 4: $\psi(t)=c|t|^\alpha$.} By~\eqref{eq:stablefe}, $\psi(n^{1/\alpha}t)=n\psi(t)$, so for integers $p,q\ge1$,
$q\,\psi((p/q)^{1/\alpha}t)=\psi(p^{1/\alpha}t)=p\psi(t)$: $\psi(r^{1/\alpha}t)=r\psi(t)$ for rational $r>0$, hence for real $r>0$ by continuity. With $t=1$ and $r=s^\alpha$,
$\psi(s)=c\,s^\alpha$ for $s>0$, where $c=\psi(1)$, and $\psi$ is even. If $c=0$ then $\psi\equiv0$; so $c>0$.

\emph{Step 5: $\alpha\le2$.} If $\alpha>2$, then $\frac{2(1-\varphi(t))}{t^2}\to0$ as $t\to0$. By Fatou's lemma and $\frac{2(1-\cos tX)}{t^2}\to X^2$,
$\mathbb EX^2\le\liminf_{t\to0}\frac{2(1-\varphi(t))}{t^2}=0$, so $X=0$, a contradiction.

\emph{Sufficiency.} For $\alpha=2$, $e^{-ct^2}$ is Gaussian. For $0<\alpha<2$, take $\Pi(\dd x)=c'|x|^{-1-\alpha}\dd x$, a L\'evy measure since $\alpha<2$ (near $0$)
and $\alpha>0$ (at infinity). By symmetry of $\Pi$ the odd part of the integrand in~\eqref{eq:lk} integrates to $0$, and with $b=\sigma^2=0$,
\[
\psi(t)=-\int(1-\cos tx)\,c'|x|^{-1-\alpha}\dd x=-c'|t|^\alpha\int(1-\cos u)|u|^{-1-\alpha}\dd u,
\]
substituting $u=tx$. Choosing $c'$ makes this $-c|t|^\alpha$, and Theorem~\ref{thm:lk}(a) says that $e^{-c|t|^\alpha}$ is a characteristic function. Finally
$\varphi(t)^n=e^{-nc|t|^\alpha}=\varphi(n^{1/\alpha}t)$: $X$ is stable with $a_n=n^{1/\alpha}$.
\end{proof}

So for $\alpha<1$ the average $S_n/n$ of $n$ copies has the law of $n^{1/\alpha-1}X$: averaging makes things \emph{worse}. For $\alpha=1$ it changes
nothing (Poisson's observation), and only for $\alpha>1$ does it help, at the slow rate $n^{1/\alpha-1}$. Figure~\ref{fig:stable} shows four of these laws. Apart
from $\alpha=1$ and $\alpha=2$ (and $\alpha=\frac12$ in the one-sided case of the next section), their densities have no closed form.

\begin{figure}[ht]
\centering\fig{ch05-stable}
\caption{Densities of the symmetric stable laws with characteristic function $e^{-|t|^\alpha}$, for $\alpha=\frac12$, $1$ (Cauchy), $\frac32$ and $2$ (the
Gaussian with variance $2$), computed by Fourier inversion. Smaller $\alpha$ means a sharper peak and heavier tails, $P(|X|>x)\approx C_\alpha x^{-\alpha}$.}
\label{fig:stable}
\end{figure}

Which i.i.d.\ sums are attracted to a stable law? For the Gaussian, a finite variance suffices (and is almost necessary). For the others, the answer is
in the tail.

\begin{theorem}[domain of attraction, symmetric case]\label{thm:domain}
Let $X$ be symmetric with $P(|X|>x)\sim Cx^{-\alpha}$ as $x\to\infty$, where $C>0$ and $0<\alpha<2$. Then
\[
\frac{X_1+\dots+X_n}{n^{1/\alpha}}\Rightarrow Y,\qquad\varphi_Y(t)=e^{-C\kappa_\alpha|t|^\alpha},\quad\kappa_\alpha=\alpha\int_0^\infty\frac{1-\cos u}{u^{1+\alpha}}\dd u.
\]
\end{theorem}
\begin{proof}
Apply Theorem~\ref{thm:arrays} to the law $\mu_n$ of $X/n^{1/\alpha}$, with $\Pi(\dd x)=\frac{C\alpha}2|x|^{-1-\alpha}\dd x$, which has no atoms. (i) By symmetry,
$n\mu_n((y,\infty))=\frac n2P(|X|>yn^{1/\alpha})\to\frac C2y^{-\alpha}=\Pi((y,\infty))$, and the same on the left. (iii) By symmetry the truncated means are $0$; take $b=0$.
(ii) with $\sigma^2=0$: for $z>0$, $\mathbb E[X^2;|X|\le z]=\int_0^z2x\,P(x<|X|\le z)\dd x\le\int_0^z2x\,P(|X|>x)\dd x$. Choose $x_0$ with $P(|X|>x)\le2Cx^{-\alpha}$ for
$x\ge x_0$; then $\mathbb E[X^2;|X|\le z]\le x_0^2+\frac{4C}{2-\alpha}z^{2-\alpha}$. With $z=\varepsilon n^{1/\alpha}$,
\[
n\int_{|x|\le\varepsilon}x^2\,\mu_n(\dd x)=n^{1-2/\alpha}\,\mathbb E[X^2;|X|\le\varepsilon n^{1/\alpha}]\le x_0^2n^{1-2/\alpha}+\frac{4C}{2-\alpha}\varepsilon^{2-\alpha},
\]
whose $\limsup$ tends to $0$ with $\varepsilon$, because $\alpha<2$. The limit exponent is $-\int(1-\cos tx)\Pi(\dd x)=-C\alpha|t|^\alpha\int_0^\infty(1-\cos u)u^{-1-\alpha}\dd u$.
\end{proof}

For the Cauchy law, $P(|X|>x)=1-\frac2\pi\arctan x\sim\frac2{\pi x}$ and $\kappa_1=\int_0^\infty\frac{1-\cos u}{u^2}\dd u=\frac\pi2$, so the theorem returns $e^{-|t|}$, as it should.
The tail condition can be weakened to $P(|X|>x)=x^{-\alpha}\ell(x)$ with $\ell$ \emph{slowly varying} ($\ell(\lambda x)/\ell(x)\to1$ for every $\lambda>0$), with
other scales $a_n$; this is necessary as well as sufficient (Gnedenko and Doeblin, around 1940). Non-symmetric stable laws exist too; their characteristic
functions are $\exp\bigl(-c|t|^\alpha(1-i\beta\operatorname{sign}(t)\tan\frac{\pi\alpha}2)+i\gamma t\bigr)$ for $\alpha\ne1$, with a skewness $\beta\in[-1,1]$ (Feller's volume~II,
Chapter~XVII). We meet one of them now.

\begin{think}
For $1<\alpha<2$ the mean exists but the variance does not. Is the sample mean of such data consistent? At what rate does it converge? How would you
report an error bar for it, and why would the usual $\pm2\hat\sigma/\sqrt n$ mislead you?
\end{think}

\section{Waiting to get ahead: the one-sided stable law of index $\frac12$}

A fair gambler waits for the first time she is one unit ahead. Volume~3 showed that she gets there almost surely, but that the expected waiting time is
infinite. What is the law of the wait? Brownian motion answers first.

\begin{lemma}\label{lem:schlomilch}
For $p,q>0$, $\displaystyle\int_0^\infty t^{-1/2}e^{-pt-q/t}\dd t=\sqrt{\frac\pi p}\,e^{-2\sqrt{pq}}$ and $\displaystyle\int_0^\infty t^{-3/2}e^{-pt-q/t}\dd t=\sqrt{\frac\pi q}\,e^{-2\sqrt{pq}}$.
\end{lemma}
\begin{proof}
Call the integrals $J_{1/2}$ and $J_{3/2}$. The substitution $t=\frac q{ps}$ exchanges $pt$ and $\frac qt$ and turns $J_{3/2}$ into $\sqrt{p/q}\,J_{1/2}$, so
$\sqrt q\,J_{3/2}=\sqrt p\,J_{1/2}$. Now $pt+\frac qt=u^2+2\sqrt{pq}$ with $u=\sqrt{pt}-\sqrt{q/t}$, which increases from $-\infty$ to $\infty$, and
$\dd u=\frac12(\sqrt pt^{-1/2}+\sqrt qt^{-3/2})\dd t$. Hence
\[
\frac12\bigl(\sqrt pJ_{1/2}+\sqrt qJ_{3/2}\bigr)=e^{-2\sqrt{pq}}\int_{-\infty}^\infty e^{-u^2}\dd u=\sqrt\pi\,e^{-2\sqrt{pq}},
\]
and the two halves are equal.
\end{proof}

\begin{theorem}[hitting times of Brownian motion]\label{thm:hitting}
Let $T_a$ be the hitting time of $a>0$ by a Brownian motion.
\begin{enumerate}[label=(\alph*)]
\item $T_a$ has the density $\dfrac a{\sqrt{2\pi}}\,t^{-3/2}e^{-a^2/(2t)}$ on $(0,\infty)$, and $P(T_1>t)\sim\sqrt{2/(\pi t)}$ as $t\to\infty$.
\item $\mathbb Ee^{-\lambda T_a}=e^{-a\sqrt{2\lambda}}$ for $\lambda\ge0$.
\item $T_1,T_2-T_1,T_3-T_2,\dots$ are i.i.d., and $T_n\overset{\mathrm d}=n^2T_1$. So $T_1$ is a one-sided stable law: a sum of $n$ copies is $n^2$ times a copy (index $\frac12$).
\end{enumerate}
\end{theorem}
\begin{proof}
(a) Differentiate $P(T_a\le t)=2(1-\Phi(a/\sqrt t))$ (Volume~3, Corollary~11.13): the derivative is $2\phi(a/\sqrt t)\cdot\frac a2t^{-3/2}$. Also
$P(T_1>t)=2\Phi(t^{-1/2})-1=2\int_0^{t^{-1/2}}\phi\sim2\phi(0)t^{-1/2}$.
(b) By (a) and Lemma~\ref{lem:schlomilch} with $p=\lambda$, $q=\frac{a^2}2$: $\mathbb Ee^{-\lambda T_a}=\frac a{\sqrt{2\pi}}\sqrt{\frac{2\pi}{a^2}}e^{-2\sqrt{\lambda a^2/2}}=e^{-a\sqrt{2\lambda}}$.
(For $\lambda=0$ both sides are $1$.)
(c) $T_k$ is a stopping time, finite almost surely, and $B_{T_k}=k$. By the strong Markov property (Theorem~\ref{thm:strongmarkov}), $W_s=B_{T_k+s}-k$ is a Brownian
motion independent of $\mathcal G_{T_k}$, and $T_{k+1}-T_k$ is its hitting time of $1$. So $T_{k+1}-T_k$ has the law of $T_1$ and is independent of
$T_1,\dots,T_k$, which are $\mathcal G_{T_k}$-measurable. By scaling (Volume~3, Proposition~11.5), $\tilde B_s=\frac1nB_{n^2s}$ is a Brownian motion and
$T_n=n^2\tilde T_1$, where $\tilde T_1$ is the hitting time of $1$ by $\tilde B$; so $T_n\overset{\mathrm d}=n^2T_1$. (In terms of (b): $(e^{-\sqrt{2\lambda}})^n=e^{-\sqrt{2n^2\lambda}}$.)
\end{proof}

The law of $T_1$ is the \emph{L\'evy distribution}. Its L\'evy measure is $\frac{\dd x}{\sqrt{2\pi}\,x^{3/2}}$ on $(0,\infty)$ (Exercise~\ref{ex:levymeasure}): a stream of
upward jumps, many small and a few enormous. Now back to the gambler.

\begin{theorem}[the waiting time of a fair gambler]\label{thm:gambler}
Let $S$ be the coin-tossing walk and $\tau_n=\min\{k:S_k=n\}$. Then $\tau_n/n^2\Rightarrow T_1$: for every $t>0$,
\[
P(\tau_n\le n^2t)\to P(T_1\le t)=2\bigl(1-\Phi(t^{-1/2})\bigr).
\]
\end{theorem}
\begin{proof}
Let $m=\lfloor n^2t\rfloor$. Then $\{\tau_n\le m\}=\{\max_{k\le m}S_k\ge n\}$. By Corollary~\ref{cor:max}, $M_m=\max_{k\le m}S_k/\sqrt m\Rightarrow|Z|$ with $Z\sim N(0,1)$,
a law with a continuous distribution function, and $n/\sqrt m\to t^{-1/2}$. So $P(M_m\ge n/\sqrt m)\to P(|Z|\ge t^{-1/2})=2(1-\Phi(t^{-1/2}))$ (for a fixed small
$\delta$, $P(M_m\ge t^{-1/2}+\delta)\le P(M_m\ge n/\sqrt m)\le P(M_m\ge t^{-1/2}-\delta)$ for large $n$, and both bounds converge).
\end{proof}

So the wait to be $n$ units ahead is about $n^2$ times a L\'evy variable: half the time it is less than $n^2/\Phi^{-1}(\frac34)^2\approx2.2\,n^2$, but its
tail $P(\tau_n>n^2t)\approx\sqrt{2/(\pi t)}$ is so heavy that the mean is infinite. Figure~\ref{fig:firstpassage} compares the exact law of $\tau_{20}$ with the limit.
The sum $\tau_n=\tau_1+(\tau_2-\tau_1)+\dots$ of i.i.d.\ waits is in the domain of attraction of a one-sided stable law: a non-symmetric instance of
Theorem~\ref{thm:arrays}, with $P(\tau_1>k)\sim\sqrt{2/(\pi k)}$.

\begin{figure}[ht]
\centering\fig{ch05-firstpassage}
\caption{The exact law of $\tau_{20}/400$, the time for the coin-tossing walk to reach $20$ divided by $20^2$, as a histogram (bins of width $0.1$; computed
from $P(\tau_n=k)=\frac nkP(S_k=n)$), against the L\'evy density $\frac1{\sqrt{2\pi}}t^{-3/2}e^{-1/(2t)}$. The mode is at $t=\frac13$; the mean is infinite.}
\label{fig:firstpassage}
\end{figure}

\section{Exercises}

\begin{exercise}[$\star$]\label{ex:idexamples}
Show that the gamma law with density $\frac{x^{s-1}e^{-x}}{\Gamma(s)}$ on $(0,\infty)$ is infinitely divisible, and that its L\'evy measure is $\frac{e^{-x}}x\dd x$ on $(0,\infty)$.
Deduce that the geometric law on $\{0,1,2,\dots\}$ (a mixture of Poisson laws with exponential mean) is infinitely divisible.
\end{exercise}

\begin{exercise}[$\star$]\label{ex:bounded}
Show that a bounded infinitely divisible variable is constant.
\end{exercise}

\begin{exercise}[$\star$]\label{ex:stablesum}
Let $X$ and $Y$ be independent, symmetric stable with characteristic functions $e^{-c_1|t|^\alpha}$ and $e^{-c_2|t|^\alpha}$. Show that $\lambda X+\mu Y$
has the law of $(c_1|\lambda|^\alpha+c_2|\mu|^\alpha)^{1/\alpha}X/c_1^{1/\alpha}$. For $\alpha<1$, compare the spread of the average of $n$ copies with that of one copy.
\end{exercise}

\begin{exercise}[$\star\star$]\label{ex:levymeasure}
Show that $\int_0^\infty(1-e^{-\lambda x})\frac{\dd x}{\sqrt{2\pi}\,x^{3/2}}=\sqrt{2\lambda}$, and deduce that the L\'evy measure of $T_1$ is $\frac{\dd x}{\sqrt{2\pi}\,x^{3/2}}$ on
$(0,\infty)$, with no Gaussian part. (Use Laplace transforms instead of characteristic functions; for positive variables they also determine the law.)
\end{exercise}

\begin{exercise}[$\star\star$]\label{ex:boundary}
Let $X$ be symmetric with $P(|X|>x)=x^{-2}$ for $x\ge1$, so that $\mathbb EX^2=\infty$ but only barely. Show that $\frac{X_1+\dots+X_n}{\sqrt{n\log n}}\Rightarrow N(0,1)$.
\end{exercise}

\begin{exercise}[$\star\star$]\label{ex:paretodomain}
Let $X$ be symmetric with $P(|X|>x)=x^{-\alpha}$ for $x\ge1$, $0<\alpha<2$. Show that $X$ is not stable, but is in the domain of attraction of a stable law,
and give its characteristic function in terms of $\kappa_\alpha$. Compute $\kappa_1$.
\end{exercise}

\begin{exercise}[$\star\star\star$]\label{ex:holtsmark}
(Holtsmark's field, one-dimensional version.) Charges sit at the points of a Poisson process of rate $\rho$ on $\R$, each with an independent random sign, and a
charge at distance $r$ exerts a force $\pm r^{-p}$ at the origin, $p>\frac12$. Let $F_R$ be the total force of the charges in $[-R,R]$. Show that
$\mathbb Ee^{itF_R}=\exp\bigl(-\rho\int_{-R}^R(1-\cos(t|x|^{-p}))\dd x\bigr)$, and that $F_R$ converges in law as $R\to\infty$ to a symmetric stable law of index $1/p$.
\end{exercise}

\section{Hints}

\begin{hint}{ex:idexamples}
The gamma$(s)$ law has characteristic function $(1-it)^{-s}$, and gamma$(s/n)^{*n}=$ gamma$(s)$. For the L\'evy measure, check that
$\int_0^\infty(1-e^{-\lambda x})\frac{e^{-x}}x\dd x=\log(1+\lambda)$ by differentiating in $\lambda$. For the geometric law, condition on the Poisson mean.
\end{hint}
\begin{hint}{ex:bounded}
If $|X|\le M$ and $X$ is a sum of $n$ i.i.d.\ pieces, each piece is bounded by $M/n$. Look at the variance.
\end{hint}
\begin{hint}{ex:stablesum}
Multiply characteristic functions.
\end{hint}
\begin{hint}{ex:levymeasure}
Integrate by parts, then use $\int_0^\infty x^{-1/2}e^{-\lambda x}\dd x=\sqrt{\pi/\lambda}$.
\end{hint}
\begin{hint}{ex:boundary}
Theorem~\ref{thm:arrays} with $a_n=\sqrt{n\log n}$: compute $\mathbb E[X^2;|X|\le z]=2\log z$ for $z\ge1$.
\end{hint}
\begin{hint}{ex:paretodomain}
$X$ never takes values in $(-1,1)$; does $X_1+X_2$? For $\kappa_1$, integrate
by parts and use $\int_0^\infty\frac{\sin u}u\dd u=\frac\pi2$.
\end{hint}
\begin{hint}{ex:holtsmark}
Given the number $N$ of points in $[-R,R]$ (Poisson of mean $2\rho R$), the points are i.i.d.\ uniform (Volume~3, Theorem~9.4).
\end{hint}

\section{Solutions}

\begin{solution}{ex:idexamples}
\idea A law built from a parameter that adds under convolution is infinitely divisible: cut the parameter into $n$ equal parts.
\writeup Let $\varphi$ be the characteristic function of the gamma$(s)$ law. Differentiating under the integral and integrating by parts,
$\varphi'(t)=\frac{is}{1-it}\varphi(t)$, so $\varphi(t)=(1-it)^{-s}$ (principal power, the solution with $\varphi(0)=1$). Hence $\varphi=\bigl((1-it)^{-s/n}\bigr)^n$, the $n$-th power
of the gamma$(s/n)$ characteristic function. For the L\'evy measure, let
\[
h(\lambda)=\int_0^\infty(1-e^{-\lambda x})\frac{e^{-x}}x\dd x.
\]
Then $h(0)=0$ and $h'(\lambda)=\int_0^\infty e^{-(1+\lambda)x}\dd x=\frac1{1+\lambda}$, so $h(\lambda)=\log(1+\lambda)$ and
\[
\mathbb Ee^{-\lambda X}=(1+\lambda)^{-s}=\exp\Bigl(-s\int_0^\infty(1-e^{-\lambda x})\frac{e^{-x}}x\dd x\Bigr),
\]
the Laplace form of the L\'evy--Khintchine formula (as in Exercise~\ref{ex:levymeasure}) with $\Pi=s\frac{e^{-x}}x\dd x$ and no drift. If $N$ is Poisson with an exponential mean $\Lambda$,
$\mathbb Ez^N=\mathbb Ee^{\Lambda(z-1)}=\frac1{1+(1-z)}=\frac1{2-z}$, the generating function of the geometric law $P(N=k)=2^{-k-1}$. Replacing $\Lambda$ by a gamma$(\frac1n)$
mean gives a law whose $n$-fold convolution is the geometric law (the means add): the geometric law is infinitely divisible.
\end{solution}

\begin{solution}{ex:bounded}
\idea A piece of a bounded variable must be very small, because $n$ pieces that are all large at once would push the sum out of bounds.
\writeup Let $|X|\le M$ and $X\overset{\mathrm d}=Y_1+\dots+Y_n$ with $Y_k$ i.i.d. If $P(Y_1>M/n)>0$ then $P(\sum Y_k>M)\ge P(Y_1>M/n)^n>0$, impossible; similarly below.
So $|Y_1|\le M/n$, $\operatorname{Var}Y_1\le M^2/n^2$ and $\operatorname{Var}X=n\operatorname{Var}Y_1\le M^2/n$ for every $n$: $\operatorname{Var}X=0$.
\end{solution}

\begin{solution}{ex:stablesum}
\idea Stable laws with the same index form a one-parameter family under convolution, the parameter being $c$, and the scale of $\lambda X$ is $|\lambda|^\alpha c$.
\writeup $\mathbb Ee^{it(\lambda X+\mu Y)}=e^{-c_1|\lambda t|^\alpha-c_2|\mu t|^\alpha}=e^{-(c_1|\lambda|^\alpha+c_2|\mu|^\alpha)|t|^\alpha}$, which is the characteristic function of $sX$ with
$c_1s^\alpha=c_1|\lambda|^\alpha+c_2|\mu|^\alpha$. For the average of $n$ copies, $c=n\cdot n^{-\alpha}c_1$, so it is $n^{1/\alpha-1}X$ in law: for $\alpha<1$ the average is
\emph{more} spread out than one copy, by a factor that grows without bound.
\end{solution}

\begin{solution}{ex:levymeasure}
\idea A positive infinitely divisible variable is a sum of positive jumps; its Laplace exponent is $\int(1-e^{-\lambda x})\Pi(\dd x)$, the Laplace version of~\eqref{eq:lk}.
\writeup Integrating by parts, $\int_0^\infty(1-e^{-\lambda x})x^{-3/2}\dd x=\bigl[-2x^{-1/2}(1-e^{-\lambda x})\bigr]_0^\infty+2\lambda\int_0^\infty x^{-1/2}e^{-\lambda x}\dd x=2\lambda\sqrt{\pi/\lambda}=2\sqrt{\pi\lambda}$
(the bracket vanishes at both ends). Dividing by $\sqrt{2\pi}$ gives $\sqrt{2\lambda}$, so $\mathbb Ee^{-\lambda T_1}=\exp\bigl(-\int(1-e^{-\lambda x})\Pi(\dd x)\bigr)$ with $\Pi=\frac{\dd x}{\sqrt{2\pi}x^{3/2}}$.
This $\Pi$ is a L\'evy measure ($\int_0^1x\cdot x^{-3/2}\dd x<\infty$, which even allows the jumps to be added without compensation), and the exponent is that of a
pure-jump law with no drift and no Gaussian part. The analogue of Theorem~\ref{thm:lk}(a) for Laplace transforms (the same proof, with compound Poisson
approximations of $\Pi$ restricted to $x>\varepsilon$) and the uniqueness of Laplace transforms of positive variables identify the law.
\end{solution}

\begin{solution}{ex:boundary}
\idea At the boundary $\alpha=2$, no single jump matters but the truncated variance grows like $\log$, which the normalisation must absorb.
\writeup For $z\ge1$, $\mathbb E[X^2;|X|\le z]=\int_1^zx^2\cdot2x^{-3}\dd x=2\log z$, since $|X|$ has density $2x^{-3}$ on $[1,\infty)$. Let $a_n=\sqrt{n\log n}$ and $\mu_n$ the law of
$X/a_n$. (i) $n\mu_n(|x|>y)=n(ya_n)^{-2}=\frac1{y^2\log n}\to0$: $\Pi=0$. (iii) The truncated means vanish by symmetry. (ii)
$n\int_{|x|\le\varepsilon}x^2\dd\mu_n=\frac n{a_n^2}2\log(\varepsilon a_n)=\frac{2\log\varepsilon+\log n+\log\log n}{\log n}\to1$ for every $\varepsilon$. So $\sigma^2=1$, and
Theorem~\ref{thm:arrays} gives the limit $N(0,1)$.
\end{solution}

\begin{solution}{ex:paretodomain}
\idea The tail hypothesis of Theorem~\ref{thm:domain} holds exactly; stability fails because a stable law has a smooth positive density, while $X$ has none
near $0$.
\writeup Theorem~\ref{thm:domain} with $C=1$ gives $S_n/n^{1/\alpha}\Rightarrow Y$ with $\varphi_Y(t)=e^{-\kappa_\alpha|t|^\alpha}$. If $X$ were stable, $X_1+X_2$ would have the law of
$a_2X$, which gives no mass to $(-a_2,a_2)$ since $P(|X|<1)=0$. But for every $\delta>0$,
$P(|X_1+X_2|<\delta)\ge P(1<X_1<1+\frac\delta2)\,P(-1-\frac\delta2<X_2<-1)>0$; take $\delta=a_2$. So $X$ is not stable. For $\kappa_1$: integrating by parts,
$\int_0^\infty\frac{1-\cos u}{u^2}\dd u=\bigl[-\frac{1-\cos u}u\bigr]_0^\infty+\int_0^\infty\frac{\sin u}u\dd u=\frac\pi2$.
\end{solution}

\begin{solution}{ex:holtsmark}
\idea Given how many charges there are, they are independent and uniform, so the characteristic function of the total is a compound Poisson one;
letting the window grow, the L\'evy measure of the force is the image of $\rho\,\dd x$ under $x\mapsto\pm|x|^{-p}$, a power law of index $1/p$.
\writeup Let $N$ be the number of points in $[-R,R]$, Poisson with mean $2\rho R$; given $N=k$, the points are i.i.d.\ uniform on $[-R,R]$ (Volume~3, Theorem~9.4),
so $F_R$ is a compound Poisson variable with jump $\varepsilon U^{-p}$, where $U=|$uniform point$|$ and $\varepsilon$ a sign. Its characteristic function is
$\exp\bigl(2\rho R(\mathbb E\cos(t|V|^{-p})-1)\bigr)$ with $V$ uniform on $[-R,R]$ (the sine part vanishes by the random sign), which is
$\exp\bigl(-\rho\int_{-R}^R(1-\cos(t|x|^{-p}))\dd x\bigr)$. As $R\to\infty$ the integral increases to $2\int_0^\infty(1-\cos(tx^{-p}))\dd x$. Substituting $u=|t|x^{-p}$,
$x=(|t|/u)^{1/p}$, this is $\frac2p|t|^{1/p}\int_0^\infty(1-\cos u)u^{-1-1/p}\dd u$, finite exactly when $0<\frac1p<2$. So $\mathbb Ee^{itF_R}\to e^{-c|t|^{1/p}}$, the
symmetric stable law of index $1/p$, by L\'evy's theorem. In three dimensions, with the inverse-square law, the same computation gives index $\frac32$: this is
the law that the Norwegian physicist Johan Holtsmark found in 1919 for the electric field in a plasma, which Chandrasekhar later used for the gravitational
field of the stars.
\end{solution}

\section*{Sources and further reading}

\begin{itemize}
\item \emph{History.} S.~M.~Stigler, ``Studies in the history of probability and statistics XXXIII: Cauchy and the witch of Agnesi: an historical note on the
Cauchy distribution'', \emph{Biometrika} 61 (1974), 375--380 (Poisson in 1824, Cauchy and Bienaym\'e in 1853). A.-L.~Cauchy, ``Sur les r\'esultats moyens
d'observations de m\^eme nature, et sur les r\'esultats les plus probables'', \emph{C.\ R.\ Acad.\ Sci.\ Paris} 37 (1853), 198--206. P.~L\'evy, \emph{Calcul des
probabilit\'es}, Gauthier-Villars, 1925 (stable laws). B.~de Finetti, ``Sulle funzioni ad incremento aleatorio'', \emph{Rend.\ Accad.\ Naz.\ Lincei} (6) 10 (1929),
163--168. A.~N.~Kolmogorov, ``Sulla forma generale di un processo stocastico omogeneo'', \emph{Rend.\ Accad.\ Naz.\ Lincei} (6) 15 (1932), 805--808. P.~L\'evy, ``Sur
les int\'egrales dont les \'el\'ements sont des variables al\'eatoires ind\'ependantes'', \emph{Ann.\ Scuola Norm.\ Sup.\ Pisa} (2) 3 (1934), 337--366. A.~Ya.~Khintchine,
``A new derivation of a formula by P.~L\'evy'' (in Russian), \emph{Bull.\ Moscow State Univ.}\ 1 (1937), 1--5. F.~Mainardi and S.~Rogosin, ``The origin of infinitely
divisible distributions: from de Finetti's problem to L\'evy--Khintchine formula'', \emph{Math.\ Methods Econ.\ Finance} 1 (2006), 37--55 (the history of that decade).
J.~Holtsmark, ``\"Uber die Verbreiterung von Spektrallinien'', \emph{Ann.\ Physik} 58 (1919), 577--630.
\item B.~V.~Gnedenko and A.~N.~Kolmogorov, \emph{Limit Distributions for Sums of Independent Random Variables}, Addison-Wesley, 1954. The complete theory of
limits of triangular arrays, and domains of attraction.
\item W.~Feller, \emph{An Introduction to Probability Theory and Its Applications}, vol.~II, 2nd ed., Wiley, 1971, Chapters~VI, IX and XVII. Stable laws, infinitely divisible laws,
and the general (non-symmetric) stable characteristic functions.
\item K.~Sato, \emph{L\'evy Processes and Infinitely Divisible Distributions}, Cambridge University Press, 1999. The modern treatise, with the processes behind the laws.
\item J.~P.~Nolan, \emph{Univariate Stable Distributions}, Springer, 2020. Computation and statistics with stable laws, including methods for simulation.
\end{itemize}
